\backhead{OBJECTIVE TYPE QUESTIONS}{khred}
\begin{exlist}
\item A Daisy log file has the extension\opt{.dai}{.dlf}{.csv}{.log}
\item The yield at 15 per cent moisture of 6 Mg DM\,ha$^{-1}$ is\opt{5.1}{6.0}{7.06}{6.9} Mg\,ha$^{-1}$
\item A residual of 20 mm in a yearly balance indicates\opt{a good closure}{a missing term}{a wet year}{a dry year}
\item A daily log over 20 years has about\opt{365}{7300}{20\,000}{175\,000} rows
\item Transpiration efficiency is expressed in\opt{kg\,m$^{-3}$}{mm}{kg\,ha$^{-1}$}{dimensionless}
\end{exlist}
\objkey{1.\ (b),\quad 2.\ (c),\quad 3.\ (b),\quad 4.\ (b),\quad 5.\ (a)}
\dykband{EXERCISES IN SWAFL}
{\small Open SWAFL (desktop or online) with the benchmark project. The chapter of the \href{https://github.com/halilunaziru73-creator/Coupled-Water-Fluxes-Pore-Scale-Dynamics-Precision-Irrigation-Mulching-Feedback-Loops/raw/main/Software/v1.1.0/SWAFL_Desktop/docs/manual_v1.1.0/SWAFL_Manual_v1.1.0.pdf}{SWAFL manual} that explains each task is given in brackets; record your values in a table and compare them with what is expected (\emph{Check}).}\par\vspace{3pt}
\begin{exlist}
\item {\color{khblue}[Manual, Chapter 30]} Open the raw output of a run and identify the header, names, units and rows.\par\noindent{\small\itshape Check: as in the figure of this chapter.}
\item {\color{khblue}[Manual, Chapter 30]} Open the efficiency indices of an experiment and compute one by hand.\par\noindent{\small\itshape Check: the hand value agrees.}
\item {\color{khblue}[Manual, Chapter 30]} Find where the results are saved and open one log file in a text editor.\par\noindent{\small\itshape Check: the folder holds every log.}
\item {\color{khblue}[Manual, Chapter 30]} Open \emph{Why a run gave no results} and explain one failure.\par\noindent{\small\itshape Check: the line of the set-up is shown.}
\end{exlist}
